% !TeX program = lualatex
% Skeleton of a method-paper reading note. Fictional placeholder content throughout.
\documentclass[a4paper,11pt]{article}
\usepackage[margin=23mm]{geometry}
\usepackage{amsmath,amssymb,graphicx,booktabs,array,listings,xcolor,hyperref}
\usepackage[most]{tcolorbox}
\lstset{basicstyle=\ttfamily\scriptsize,breaklines=true,numbers=left,numberstyle=\tiny,frame=single}
% Shared macros for a method-paper reading note (skeleton; fictional defaults).
% Set these once per note. NoteDir is the note's directory relative to the workspace root.
\newcommand{\NoteDir}{reading/example-note}
\newcommand{\PDFURL}{file:///path/to/the/stored/paper.pdf}
% Labels (translate here, once).
\newcommand{\LabelOrig}{Original}
\newcommand{\LabelTrans}{Faithful translation}
\newcommand{\LabelSrc}{Source}
\newcommand{\LabelAdded}{Added by this note}
\newcommand{\LabelCard}{Background card}
\newcommand{\LabelFrom}{Taken from earlier units}
\newcommand{\LabelTo}{Handed to later units}
% \srcp{journal page}{pdf page}{label}: link to a page of the stored PDF
\newcommand{\srcp}[3]{\href{\PDFURL\#page=#2}{#3, p.#1}}
% provenance tags (R5a): source, and material the note adds (not evidence for a claim)
\newcommand{\src}[1]{{\small\color{teal}[\LabelSrc: #1]}}
\newcommand{\mine}[1]{{\small\color{brown}[\LabelAdded: #1]}}
% role heading: role | concrete point
\newcommand{\role}[2]{\subsubsection{#1 | #2}}
% term registration for the index
\newcommand{\defterm}[1]{\phantomsection\label{def:#1}}
% named-work heading above an original/translation pair
\newcommand{\workhead}[2]{\par\medskip\noindent{\small\bfseries\textcolor{blue}{#1}\ --\ #2}\par\nopagebreak}
\newenvironment{orig}{\par\smallskip\noindent\textbf{\LabelOrig}\par\begin{quote}\small\sffamily}{\end{quote}}
\newenvironment{trans}{\par\smallskip\noindent\textbf{\LabelTrans}\par\begin{quote}\small}{\end{quote}}
% \origin{name}: input NoteDir/orig/name.tex (extracted source text kept in its own file)
\newcommand{\origin}[1]{\begin{orig}\input{\NoteDir/orig/#1.tex}\end{orig}}
% background card
\newenvironment{bgcard}{\par\medskip\begin{tcolorbox}[colback=black!3,colframe=teal!60!black,title={\LabelCard},fonttitle=\bfseries\small,boxrule=0.5pt,left=2mm,right=2mm,top=1mm,bottom=1mm,breakable]\small}{\end{tcolorbox}}
\newcommand{\cardrow}[2]{\par\noindent\textbf{#1}: #2}
% \paperfig{file}{width}{journal page}{pdf page}{reading of the figure}: a figure cut out of the source
\newcommand{\paperfig}[5]{%
  \begin{figure}[htbp]\centering
  \includegraphics[width=#2\linewidth]{\NoteDir/figures/#1}
  \caption{Cut out of the source (\srcp{#3}{#4}{page}), unaltered. #5}
  \end{figure}}
% authors' code: \authorcode{file}{first}{last}{caption}, then \codeid{file}
\newcommand{\authorcode}[4]{\lstinputlisting[language=Python,firstline=#2,lastline=#3,firstnumber=#2,caption={#4}]{\NoteDir/code/#1}}
% notes under an author-code excerpt: what each block does / where it comes from / what differs and what that changes
% needs \usepackage{needspace} so the heading is not left alone at the foot of a page
\newcommand{\codenotes}[1]{\Needspace{12\baselineskip}\par\smallskip\noindent{\small\bfseries Code notes}\par\nopagebreak
  {\footnotesize\begin{longtable}{p{1.7cm}p{4.3cm}p{3.6cm}p{4.9cm}}\toprule Lines & Does & From & Difference and effect\\\midrule\endhead #1\\\bottomrule\end{longtable}}}
\newcommand{\codeid}[1]{\par{\small\textbf{Code source:} AUTHORS, REPOSITORY, commit \texttt{HASH}, \nolinkurl{#1}; licence: LICENCE. State whether this is the version used in the paper.}\par}
% unit chain box: what this unit takes from earlier units and what it hands on
\newcommand{\chain}[2]{\par\medskip\noindent\fbox{\parbox{0.96\linewidth}{\small\textbf{\LabelFrom:} #1\par\textbf{\LabelTo:} #2}}\par\medskip}

\begin{document}
\title{Reading note: ``A Fictional Ordering Method'' (Doe, 2001)}
\date{}
\maketitle
\tableofcontents

\section{Scope, reader floor and spine}
\subsection{The paper and what was read}
Identity, version, DOI and stored copy; what was read and what was not (supplement, code,
data); what is assumed of the reader (mathematics, domain). \mine{Placeholder.}

\subsection{Spine and section map}
The question, the claims in dependency order, and for each unit its parent, what it needs
and what it produces (kept as data in \texttt{coverage\_spec.json}).

\subsection{Lineage overview}\label{sec:lineage}
One row per earlier work (the paper's inherited, compared, criticised ones): made for; the
weakness the paper names; what the method keeps or changes; what checking the source found;
the section that treats it. Reference lineages the code does not show are marked; later
developments are a separate auxiliary section. \mine{Placeholder.}

\subsection{Values}\label{sec:values}
One table: the code's argument (the paper's symbol), the paper's value, the code's default and
the rule that sets it, the values used in each run of this note and of the companion.
\mine{Placeholder.}

\section{Section 2 of the paper: the method}
\chain{an ordering problem stated in the introduction}{the graph construction}
\role{Background}{what the problem needs}
\begin{bgcard}
\cardrow{Problem}{items without a reference point must be ordered.}
\cardrow{Why the previous tool is not enough}{state a concrete failure of the earlier method.}
\cardrow{Idea}{state it in words before the definition.}
\end{bgcard}
\role{Mathematics and implementation}{the step, its reason, the authors' code}
Correspondence table between the paper's step and the lines of the authors' code. After each
excerpt, \verb|\codenotes{...}|: per block, what it does in plain words, where it comes from,
and what differs and what that changes.
\role{Where the figure's difference comes from}{steps and lines, marked as the paper's statement or this note's reading}
For each main-text figure that compares methods or conditions: the blocks of the code that make
the difference.
\role{Original and translation}{the passage as published}
\workhead{Doe (2001)}{Sec.~2, p.\,3}
% \origin{sec2}  % extracted source text in orig/sec2.tex, then the faithful translation:
\begin{trans}Faithful translation of the passage goes here.\end{trans}

\section{Supplement items the main text relies on}
Each code, procedure, claim-bearing figure, derivation, data or validation item, read at the
place where the main text points to it.

\section{Execution and reproduction}
Environment and versions, runs labelled by provenance, what was reproduced and what was not.

\section{Status}
Resources not obtained or not read, the places tried, what is wanted, free alternatives, and any
stand-in made with the user's per-resource permission (labelled, shown beside the open entry).
\end{document}
